\documentclass[10pt]{article}
\usepackage[a4paper,top=1.55cm,bottom=1.7cm,left=1.55cm,right=1.55cm,columnsep=0.55cm]{geometry}
\usepackage{fontspec}
\usepackage{amsmath,amssymb}
\usepackage{booktabs,array}
\usepackage{multicol}
\usepackage{microtype}
\usepackage[hidelinks]{hyperref}
\setmainfont{STSong}[BoldFont=Heiti SC]
\setsansfont{Arial}
\setmonofont{Menlo}
\XeTeXlinebreaklocale "zh"
\XeTeXlinebreakskip=0pt plus 1pt
\setlength{\parindent}{1.5em}
\setlength{\parskip}{0pt}
\setlength{\abovedisplayskip}{4pt}
\setlength{\belowdisplayskip}{4pt}
\setlength{\textfloatsep}{7pt}
\setlength{\floatsep}{6pt}
\pagestyle{plain}
\newcommand{\code}[1]{\texttt{#1}}
\newcounter{algo}
\begin{document}

\setcounter{section}{2}
\section{方法}

\begin{center}
\fbox{\begin{minipage}{0.94\textwidth}\centering
\textbf{长期交互历史} $\rightarrow$
\textbf{L0 原始经历} $\rightarrow$
\textbf{L1 风险记忆卡} $\rightarrow$
\textbf{L2 风险规则}\\[3pt]
$\downarrow$ 角色条件化风险检索与证据预算\\[3pt]
\textbf{Task Advocate} $\parallel$ \textbf{Risk Critic}
$\rightarrow$ \textbf{Safety Decision Agent}
$\rightarrow$ \textbf{确定性约束优化器}
$\rightarrow$ \textbf{Safety Guard}
$\rightarrow$ \textbf{执行}
\end{minipage}}
\refstepcounter{figure}\label{fig:overview}
\vspace{2pt}
{\small 图~\thefigure：方法概览。分层记忆压缩长期历史并保留来源链；不同角色接收不同风险证据；确定性模块保留执行权限。}
\end{center}

\begin{multicols}{2}

\subsection{系统概览}

本文研究室内移动机器人在部分可观测环境中的高层语义动作选择。时刻 $t$ 的规划输入记为
\begin{equation}
\mathcal{X}_t=(g_t,\,r_t,\,e_t,\,\mathcal{A}_t),
\end{equation}
其中 $g_t$ 为任务目标及进度，$r_t$ 为机器人状态，$e_t$ 为环境观测，$\mathcal{A}_t$ 为当前允许的有限动作集合。系统维护长期交互历史 $\mathcal{H}_t$，并在有界模型上下文内提供与当前任务相关且可追溯的风险证据。

图~\ref{fig:overview} 给出整体流程。系统将执行记录存为 L0 原始经历，将其压缩为 L1 Risk Memory Card，并将跨经历的稳定模式抽象为 L2 Risk Rule。检索器为 Task Advocate、Risk Critic 和 Safety Decision Agent 分别构造证据包。三个 Agent 输出结构化建议及其引用的经历和规则 ID。确定性约束优化器据此选择动作，再由独立 Safety Guard 复核。

令三个角色的报告集合为 $\mathcal{D}_t$，匹配的 L2 规则为 $\mathcal{Q}_t$。优化器首先排除违反硬约束或 L2 禁止规则的动作，再求解
\begin{equation}
a_t^*=\arg\max_{a\in\mathcal{F}_t} U(a),\quad
\text{s.t. } R(a)\leq\rho,\; Z(a)\leq\eta,
\end{equation}
其中 $\mathcal{F}_t\subseteq\mathcal{A}_t$ 为可行动作集合，$U$ 为任务收益，$R$ 为场景与记忆联合风险，$Z$ 为观测置信度、记忆覆盖和 Agent 分歧构成的不确定性。$R$ 与 $Z$ 均为相对分数，而非校准概率。L2 规则冲突、Agent 输出无效、可行集为空或 Safety Guard 拒绝动作时，系统执行 \code{safe\_stop}、\code{observe\_again} 或请求人工复核。

\subsection{L0--L1--L2 分层风险记忆}

分层记忆统一表示风险经历、压缩长期历史，并保留抽象规则到原始案例的证据链。记忆图记为
\begin{equation}
\mathcal{G}_t=(\mathcal{E}_t,\mathcal{M}_t,\mathcal{Q}_t,\mathcal{L}_t),
\end{equation}
其中 $\mathcal{E}_t$、$\mathcal{M}_t$ 和 $\mathcal{Q}_t$ 分别表示 L0、L1 和 L2 节点，$\mathcal{L}_t$ 保存稳定 ID 之间的来源链接。

\subsubsection{L0：具身风险经历表示}

L0 保存能够支持事后核查和重新抽象的完整交互证据。第 $i$ 条经历写为
\begin{equation}
E_i=(G_i,S_i,A_i,O_i,F_i,K_i,D_i,P_i),
\end{equation}
其中 $G_i$ 为任务目标，$S_i$ 为机器人与环境状态，$A_i$ 为执行动作，$O_i$ 为动作结果，$F_i$ 为失败原因，$K_i$ 为风险类型与严重度，$D_i$ 为决策记录，$P_i$ 为轨迹来源及时间。该结构完整关联任务、状态、动作、结果和风险原因，以支持案例重建。

每条 L0 经历具有不可变的 \code{experience\_id}，并可链接传感日志、控制记录或人工标注。写入时校验必需字段、结果状态、严重度范围和 ID 唯一性。Agent 接收压缩后的 L1 表示，L0 则保存用于压缩与审计的原始证据。

\subsubsection{L1：风险记忆卡}

L1 通过确定性压缩映射 $C$ 将完整经历转化为紧凑风险记忆卡：
\begin{equation}
M_i=C(E_i)=(T_i,H_i,A_i,O_i,F_i,B_i,Q_i,P_i),
\end{equation}
其中 $T_i$ 为触发条件，$H_i$ 为危险类型与严重度，$B_i$ 为缓解措施和停止条件，$Q_i$ 为支持次数、可靠性及最近验证时间，$P_i$ 为来源信息。$T_i$ 仅保留通行余量、障碍物距离、观测置信度和电量等风险相关状态。

每张卡具有稳定的 \code{memory\_id}，并保存其 \code{experience\_id} 和来源片段。Agent 使用卡片推理，审计者可沿 $M_i\rightarrow E_i$ 回查原始轨迹。卡片同时记录检索得分分解和匹配原因，使检索结果可复现。

\subsubsection{L2：风险规则与认知演化}

L2 将多个相关 L1 卡片中的稳定风险模式表达为条件规则
\begin{equation}
Q_j=(C_j,K_j,A_j^{+},A_j^{-},B_j,N_j,c_j,z_j),
\end{equation}
其中 $C_j$ 为触发条件，$K_j$ 为风险类型，$A_j^{+}$ 与 $A_j^{-}$ 分别为建议和禁止动作，$B_j$ 为缓解措施，$N_j$ 为来源记忆 ID 集合，$c_j$ 为置信度，$z_j\in\{\text{candidate},\text{active},\text{retired}\}$ 为生命周期状态。规则仅在触发条件满足时参与当前决策；其来源集合提供 $Q_j\rightarrow M_i\rightarrow E_i$ 的完整追溯路径。

新轨迹形成 L0 经历，相似经历更新 L1 的支持统计与可靠性，重复且充分支持的模式形成或强化候选 L2 规则。矛盾证据降低规则置信度、修订触发范围或使规则停用。最低支持阈值避免单次失败直接生成强规则。

运行时仅匹配 \textit{active} 规则。若命中规则存在共同且未被禁止的建议动作，系统保留其交集；否则，确定性冲突闸门停止执行并请求人工复核。

为使长期更新可复现，每个风险模式维护统计量 $n_j=(n_j^{+},n_j^{-},n_j^{0})$，分别表示缓解成功、风险失败和人工中止的支持次数。缺少显式质量标注时，规则可靠性可由带平滑的证据估计给出：
\begin{equation}
c_j=\frac{n_j^{+}+n_j^{-}+\alpha}
{n_j^{+}+n_j^{-}+n_j^{0}+\alpha+\beta},
\end{equation}
其中 $\alpha,\beta>0$ 抑制小样本产生的极端置信度。候选规则在满足最低来源数、跨场景一致性和置信度阈值后转为 \textit{active}；持续冲突、长期未验证或被反例否定时降为 \textit{candidate} 或 \textit{retired}。每次状态变化均创建新版本，并保留旧版本和触发更新的来源 ID。

失败案例用于识别禁止动作和停止条件，成功案例用于验证缓解措施。规则仅在两类证据均被覆盖时扩展到相似场景，以限制失败记忆累积导致的过度保守。

\subsection{角色条件化的风险感知检索}

给定输入 $\mathcal{X}_t$ 和角色 $r\in\{\mathrm{adv},\mathrm{crit},\mathrm{dec}\}$，检索器根据机器人类型、障碍物状态、通行余量、障碍物距离、观测置信度和任务目标计算 L0 经历的确定性场景相似度。正相似度候选的角色得分为
\begin{equation}
\label{eq:role-score-zh}
s_i^{(r)}=\sum_{k\in\mathcal{K}}w_k^{(r)}\phi_{ik},\quad
\mathcal{K}=\{\mathrm{sim},\mathrm{sev},\mathrm{rel},\mathrm{rec},\mathrm{fit}\},
\end{equation}
其中 $\phi_{ik}\in[0,1]$ 表示场景相似度、风险严重度、证据可靠性、时效性或角色适配度，角色权重 $w_k^{(r)}$ 之和为 1。可靠性优先采用显式质量评分，否则采用平滑后验均值；时效性随验证间隔衰减。

角色条件化由 $w_k^{(r)}$ 与 $\mathrm{fit}$ 共同实现。Advocate 强调成功经历和有效缓解措施，Critic 强调失败、高严重度风险和停止条件，Decision 平衡正反案例。因此，三个角色在推理前获得不同的经验分布。

当前原型采用如下可审计基线权重，分量顺序为 $(\mathrm{sim},\mathrm{sev},\mathrm{rel},\mathrm{rec},\mathrm{fit})$：
\begin{align}
w^{(\mathrm{adv})}&=(.35,.05,.20,.05,.35),\\
w^{(\mathrm{crit})}&=(.30,.25,.15,.05,.25),\\
w^{(\mathrm{dec})}&=(.35,.20,.20,.10,.15).
\end{align}
这些手工设定的权重构成可复现基线。消融和敏感性分析分别检验角色适配度、严重度和可靠性的作用。

为避免 Top-$k$ 被重复案例占满，系统采用贪心多样性选择。在已选集合 $S$ 下，下一条记忆为
\begin{equation}
\label{eq:diversity-zh}
i^*=\arg\max_{i\notin S}\left[s_i^{(r)}-\lambda
\max_{j\in S} d(i,j)\right],
\end{equation}
其中 $d(i,j)$ 根据风险类型、结果状态和缓解措施计算冗余度，$\lambda$ 控制相关性与多样性的权衡。候选被转换为 L1 卡片，并装入预算为 $B_r$ 的证据包：
\begin{equation}
\label{eq:budget-zh}
\sum_{M_i\in\mathcal{B}_r}\widehat{\tau}(M_i)\leq B_r,
\end{equation}
其中 $\widehat{\tau}$ 是针对中英文 JSON 的保守 Token 估计。超过预算的卡片被跳过，从而使历史规模增长时 Agent 输入仍保持有界。

证据包包含角色、检索版本、预算使用量、L1 卡片、得分分解、匹配原因和来源 ID，并单独附加命中的 L2 规则。Agent 必须在 \code{cited\_experience\_ids} 与 \code{cited\_rule\_ids} 中声明所用证据；系统拒绝输入包中不存在的 ID。该记录支持审计检索内容、检索原因和实际引用。

算法~\ref{alg:retrieval} 概括检索过程。对大小为 $N$ 的经验库，每个角色最多选取 $K$ 张卡片；特征计算为 $O(N)$，贪心去冗余为 $O(NK)$。由于 $K$ 由上下文预算固定，在线开销随 $N$ 近似线性增长。三个角色共享场景特征，但分别排序和装包。

\begin{center}
\fbox{\begin{minipage}{0.91\columnwidth}
\refstepcounter{algo}\label{alg:retrieval}
\textbf{算法~\thealgo：角色条件化风险检索}\\[-2pt]
\textbf{输入：} 当前状态 $\mathcal{X}_t$，角色 $r$，L0 库 $\mathcal{E}$，预算 $B_r$，上限 $K$。\\
1) 校验角色与输入结构；\\
2) 计算每条经历的透明场景相似度；\\
3) 计算严重度、可靠性、时效性和角色适配度；\\
4) 使用 $w^{(r)}$ 计算角色得分 $s_i^{(r)}$；\\
5) 以式~\eqref{eq:diversity-zh} 的冗余惩罚贪心选择候选；\\
6) 将候选确定性压缩为 L1 卡片；\\
7) 按式~\eqref{eq:budget-zh} 装入 Token 预算并附加匹配 L2 规则；\\
8) 返回含得分分解、匹配原因和来源 ID 的证据包。
\end{minipage}}
\end{center}

\end{multicols}
\end{document}
